%!TEX root = ../../main.tex
%% ===============================================================
%% 10.1 정보 입도 한계
%% 파일: chapters/10_discussion/01_granularity.tex
%% 상위: chapters/10_discussion/chapter.tex
%% 정의 라벨: sec:disc-granularity
%% 참조 라벨: ch:results
%% 포함 객체: -
%% 인용: katona2019time
%% 상태: 초안 (docs/OUTLINE.md 참고)
%% ===============================================================

\section{정보 입도 한계}
\label{sec:disc-granularity}

제 \ref{ch:results} 장의 결과는 하나의 해석으로 수렴한다. 분 단위 공개 텔레메트리는 교전 결과의 \emph{거시적} 결정 요인(누적 자원 격차, 오브젝트 통제, 생존 인원)을 담고 있지만 \emph{미시적} 결정 요인(스킬 적중, 진입 타이밍, 쿨다운)은 담고 있지 않다. 거시 요인이 결과를 지배하는 후반·일방 교전에서는 AUC가 $.8$에 이르고, 미시 요인이 결과를 가르는 대등한 교전에서는 우연 수준에 가깝다. 이는 관측 정보의 상호정보량 관점에서 자연스럽다. 라벨 $Y$와 관측 $X$의 상호정보량 $I(X;Y)$는 관측 해상도가 거칠수록 작아지며, 어떤 학습기도 $I(X;Y)$가 허용하는 이상의 판별력을 얻을 수 없다. 리플레이 기반의 초 단위 관측을 전제한 선행 연구 \cite{katona2019time}가 높은 정확도를 보고한 것과 본 결과의 대비는 바로 이 해상도 차이로 설명된다.
